% Part: lambda-calculus
% Chapter: introduction
% Section: unique-readability

\documentclass[../../../include/open-logic-section]{subfiles}

\begin{document}

\olfileid{lam}{syn}{unq}
\olsection{يکتا لوستونتيا}

کېدے شي وپوښتو چې ايا د هر ترم د جوړېدو يوازې يوه لاره شته؛ هو،
همداسې ده۔ د هر لامبډا ترم د جوړولو او تفسير يوازې يوه لاره شته۔
په لاندې بحث کښې \emph{جوړونه} هغه کړنلاره ده چې د
\olref[trm]{defn:term} د جوړولو قاعدې يو يا څو ځله کاروي او ترم
ترې جوړوي۔

\begin{lem}\ollabel{lem:term-start}
ترم يا په متغير پيلېږي يا په قوس۔
\end{lem}

\begin{proof}
يو څيز هله ترم ګڼل کېږي چې د \olref[trm]{defn:term} له مخې جوړ
شوے وي۔ که د \olref[trm]{defn:term-var} پايله وي، نو بايد متغير
وي۔ که د \olref[trm]{defn:term-abs} يا
\olref[trm]{defn:term-app} پايله وي، نو په قوس پيلېږي۔
\end{proof}

\begin{lem}\ollabel{lem:app-start}
د تطبيق پايله يا په دوو قوسونو پيلېږي يا په يوه قوس او بيا متغير۔
\end{lem}

\begin{proof}
که $M$ د تطبيق پايله وي، نو بڼه ئې $(PQ)$ ده؛ ځکه په قوس پيلېږي۔
څرنګه چې $P$ ترم دے، د \olref{lem:term-start} له مخې يا په قوس
پيلېږي يا په متغير۔
\end{proof}

\begin{lem}\ollabel{lem:initial}
د ترم هېڅ خاصه پيلنۍ برخه پخپله ترم نۀ ده۔
\end{lem}

\begin{prob}
\olref[lam][syn][unq]{lem:initial} د ترمونو پر اوږدوالي په استقرا
ثابت کړئ۔
\end{prob}

\begin{prop}[يکتا لوستونتيا] \ollabel{prop:unq}
د هر ترم جوړونه يکتا ده۔ په بله وينا، که ترم $M$ د يوې جوړونې په
وسيله جوړ شي، نو د همدې ترم د جوړولو بله جوړونه نشته۔
\end{prop}

\begin{proof}
  دا د ترمونو پر جوړونه په استقرا ثابتوو۔
  \begin{enumerate}
    \item $M$ د $x$ په بڼه دے، چې $x$ يو متغير دے۔ د تجريد او
      تطبيق پايلې تل په قوسونو پيلېږي، نو د $M$ د جوړولو دپاره نۀ شي
      کارېدے؛ ځکه د $M$ جوړونه بايد د
      \olref[trm]{defn:term}\olref[trm]{defn:term-var} يوازې يو ګام
      وي۔
    \item $M$ د $(\lambd[x][N])$ په بڼه دے، چې $x$ متغير او $N$
      ترم دے۔ د \olref[trm]{defn:term}\olref[trm]{defn:term-var}
      له مخې نۀ شي جوړېدے، ځکه يوازې يو متغير نۀ دے۔ د
      \olref{lem:app-start} له مخې د تطبيق پايله هم نۀ ده۔ نو $M$
      يوازې پر~$N$ د تجريد پايله کېدے شي۔ د استقرا د فرضيې له مخې
      د $N$ جوړونه پخپله يکتا ده۔
    \item $M$ د $(PQ)$ په بڼه دے، چې $P$ او $Q$ ترمونه دي۔ څرنګه
      چې په قوس پيلېږي، د
      \olref[trm]{defn:term}\olref[trm]{defn:term-var} له مخې هم
      نۀ شي جوړېدے۔ د \olref{lem:term-start} له مخې $P$ په
      $\lambd$ نۀ شي پيلېدے، نو $(PQ)$ د تجريد پايله نۀ شي کېدے۔
      اوس فرض کړئ چې $M$ د تطبيق له لارې د جوړېدو بله طريقه هم لري؛
      يعنې د $(P'Q')$ په بڼه هم دے۔ بيا $P$ د $P'$ خاصه پيلنۍ برخه
      ده، يا برعکس؛ خو دا د \olref{lem:initial} له مخې ناشوني دي۔
      نو $P$ او $Q$ په يکتا توګه ټاکل شوي دي، او د استقرا د فرضيې
      له مخې د $P$ او $Q$ جوړونې يکتا دي۔
  \end{enumerate}
\end{proof}

د پاسنۍ قضيې د لوستلو آسانه بڼه دا ده:
\begin{prop}
  ترم $M$ يوازې له لاندې بڼو څخه يوه لرلے شي:
  \begin{enumerate}
    \item $x$، چې $x$ د $M$ له مخې په يکتا توګه ټاکل شوے متغير دے۔
    \item $(\lambd[x][N])$، چې $x$ متغير او $N$ بل ترم دے، او
      دواړه د $M$ له مخې په يکتا توګه ټاکل شوي دي۔
    \item $(PQ)$، چې $P$ او $Q$ د $M$ له مخې په يکتا توګه ټاکل شوي
      دوه ترمونه دي۔
  \end{enumerate}
\end{prop}

\end{document}
